\section{Measuring Task-Dependent Signatures}
\label{sec:methods}
\begin{table}[t]
\centering\small
\begin{tabularx}{\columnwidth}{@{}lX@{}}
\toprule Study & Population and purpose\\\midrule
Atlas & Seven models, eight benchmark conditions; prevalence and timing.\\
Matched demands & Five models, 80 math and 80 GPQA items; three activities and answer-option control.\\
Interventions & Qwen3.5-9B; five domain contrasts, 199 math and 200 items each for GPQA/planning. A2, SEAL, and historical biases are secondary comparisons.\\
\bottomrule
\end{tabularx}
\caption{Study overview. Three matched-demand models also appear in the atlas; two additional models are listed in Table~\ref{tab:models}. Intervention results do not establish model-family generality.}
\label{tab:overview}
\end{table}
\subsection{Atlas and annotations}
The primary analysis contains 21,803 traces with valid sentence annotations from seven models: DeepSeek-R1-Distill-Llama-8B; Gemma4-E4B-it, Gemma4-12B-it, and Gemma4-31B-it; and Qwen3.5-4B, Qwen3.5-9B, and Qwen3.5-27B. The larger generation release contains 27,904 traces on 3,488 prompts, including a Llama-3.1-8B-Instruct answer-only anchor excluded from the reasoning comparison. Table~\ref{tab:benchmarks} details the eight benchmark domains. Most have 500 prompts; code has 175 and safety has 313. The configuration uses one generation seed per prompt/model, temperature 0.6, top-$p$ 0.95, and top-$k$ 20. The reasoning generation cap is 65,536 tokens, subject to available context.

We segment reasoning with the repository's ThinkARM-derived splitter \citep{li2025thinkarm}. Gemma4-31B-it assigns eight binary labels: question--answering, perspective shift, perspective conflict, reconciliation, verification, backtracking, subgoal setting, and backward chaining. Labels may co-occur. The first four form the conversational group; the latter four form the cognitive-operation group. These names organize analyses, rather than asserting a psychological partition.

The taxonomies follow \citet{kim2026society} and \citet{gandhi2025habits}. Prior work reports human agreement for both families; our implementation adapts them to indexed sentence batches. Annotations are obtained with preceding reasoning context and with isolated segments. The primary features use contextual labels, retaining segments whose required annotations parse in both conditions. Context agreement measures sensitivity to context. Of 24,416 reasoning attempts, 21,803 enter the atlas, including 46 unfinished attempts. Usable-annotation coverage ranges from 67.3\% to 99.9\% across models. Completion and annotation inclusion are separate quantities (Appendix~\ref{app:audit}).

\subsection{Amount, timing, and domain identification}
For trace $i$, operation $b$, and valid segments $\mathcal V_i$, amount is
\begin{equation}
a_{ib}=|\mathcal V_i|^{-1}\sum_{s\in\mathcal V_i} y_{isb}.
\end{equation}
We separately compute prevalence in early, middle, and late thirds of normalized sentence position. Dividing each phase prevalence by their sum gives shares $q_{ibp}$; all shares are zero if the operation is absent. The timing features are $q_{ib,\mathrm{middle}}-q_{ib,\mathrm{early}}$ and $q_{ib,\mathrm{late}}-q_{ib,\mathrm{early}}$.

Within each model and annotation group, ridge multinomial logistic regression predicts the eight benchmark conditions using five prompt-disjoint folds. Nested specifications contain log token/segment counts, then four amounts, then eight timing contrasts. Preprocessing uses training folds only. We report held-out log-score gains in bits per trace, with 95\% bootstrap intervals and Benjamini--Hochberg adjustment. Support requires a positive interval lower bound and $q<.05$. Descriptive means weight traces equally within model/domain and then models equally. Similarity uses within-model RMS differences across eight rates, averaged over models. Temporal figures use ten position bins and show individual-model curves; these are marginal prevalences, not transition probabilities.

\subsection{Matched task-demand study}
We hold 80 math and 80 GPQA problems fixed while changing instructions to solve, generate approaches, or evaluate an answer. A free-response versus multiple-choice comparison tests format. The displayed contrast compares generating approaches with solving, alongside repeated-generation noise references. It includes Qwen3.5-9B, Gemma4-12B-it, DeepSeek-R1-Distill-Llama-8B, DeepSeek-R1-Distill-Qwen-7B, and OpenThinker3-7B. Rates use the first 8,192 reasoning tokens. Profile distance is the mean absolute log$_2$ ratio of smoothed rates ($\epsilon=.005$). Item-bootstrap intervals describe matched contrasts. Because the noise-adjusted demand and format quantities use different repeat references, their raw paired difference is reported separately.
